% !TEX program = xelatex
% 编译：xelatex manual.tex
\documentclass[UTF8,a4paper,11pt]{ctexart}

\usepackage[a4paper,margin=2.4cm]{geometry}
\usepackage{xcolor}
\usepackage{booktabs}
\usepackage{tabularx}
\usepackage{enumitem}
\usepackage{listings}
\usepackage{hyperref}
\usepackage{url}
\usepackage{fancyhdr}

\hypersetup{
  colorlinks=true,
  linkcolor=black,
  urlcolor=blue!50!black,
  citecolor=black,
  pdftitle={qoj-submit 使用说明},
  pdfauthor={qoj-submit}
}

\definecolor{codebg}{RGB}{247,247,248}
\definecolor{codeframe}{RGB}{210,210,214}
\definecolor{cmd}{RGB}{32,32,32}

\lstset{
  basicstyle=\small\ttfamily,
  backgroundcolor=\color{codebg},
  frame=single,
  rulecolor=\color{codeframe},
  breaklines=true,
  breakatwhitespace=false,
  columns=fullflexible,
  keepspaces=true,
  showstringspaces=false,
  tabsize=4,
  xleftmargin=0.6em,
  xrightmargin=0.6em,
  aboveskip=0.7em,
  belowskip=0.7em,
  framesep=6pt
}

\pagestyle{fancy}
\fancyhf{}
\fancyhead[L]{\small qoj-submit}
\fancyhead[R]{\small 使用说明}
\fancyfoot[C]{\thepage}
\renewcommand{\headrulewidth}{0.4pt}

\title{\textbf{qoj-submit} 使用说明}
\author{QOJ VP 的 XCPC 风格命令行提交插件}
\date{版本 1.0.0}

\setlist[itemize]{leftmargin=1.6em,itemsep=0.25em,topsep=0.4em}
\setlist[enumerate]{leftmargin=1.8em,itemsep=0.25em,topsep=0.4em}
\setlist[description]{leftmargin=1.2em,itemsep=0.35em,topsep=0.4em}

\newcommand{\cmd}[1]{\texttt{\detokenize{#1}}}
\newcommand{\opt}[1]{\texttt{\detokenize{#1}}}

\begin{document}
\maketitle
\tableofcontents
\newpage

\section{简介}

\cmd{qoj-submit} 在 Linux 下模拟 XCPC 现场的提交方式，用于 QOJ（\url{https://qoj.ac}）的 Virtual Participation（VP）。绑定比赛后，在工作目录中：

\begin{lstlisting}
./submit a.cpp      # 提交 A 题
./submit b.cpp      # 提交 B 题
./submit c.py       # 提交 C 题
\end{lstlisting}

插件是单个可执行 Python 脚本，只依赖 Python~3.8+ 标准库，无需 \cmd{pip} 安装第三方包。

\section{环境要求}

\begin{itemize}
  \item 操作系统：Linux（设计目标）；Windows / macOS 上只要有 Python~3 也可运行。
  \item Python~3.8 或更高版本。
  \item 能访问 QOJ 的网络。站点有 Cloudflare 防护，登录方式见第~\ref{sec:login}~节。
\end{itemize}

\section{安装}

将仓库中的 \cmd{submit} 拷到 Linux 后赋予可执行权限：

\begin{lstlisting}
chmod +x submit
\end{lstlisting}

之后有两种用法：

\begin{enumerate}
  \item 把 \cmd{submit} 放到 \texttt{\$PATH} 中（例如 \texttt{\textasciitilde/bin}），任意目录都可调用。
  \item 每场比赛执行 \cmd{init} 时，脚本会自动复制一份到当前目录，之后直接使用 \cmd{./submit}。
\end{enumerate}

查看帮助与版本：

\begin{lstlisting}
./submit --help
./submit --version
\end{lstlisting}

\section{登录}
\label{sec:login}

账号信息保存在本机配置文件中，每台机器一般只需登录一次。QOJ 有 Cloudflare，\textbf{推荐粘贴浏览器 Cookie}，不推荐明文密码。

\subsection{从浏览器复制 Cookie}

\begin{enumerate}
  \item 用浏览器登录 \url{https://qoj.ac}。
  \item 按 \cmd{F12} 打开开发者工具，切换到 \emph{Network}（网络）。
  \item 刷新页面或点击任意请求，在 Request Headers 中复制完整的 \cmd{cookie}。
  \item 若随后提示被 Cloudflare 拦截，再复制同一请求中的 \cmd{user-agent}。Cookie 与 User-Agent 必须来自\textbf{同一次浏览器会话}，通常需要包含 \cmd{cf\_clearance}。
\end{enumerate}

交互式登录：

\begin{lstlisting}
./submit login
\end{lstlisting}

按提示粘贴 Cookie；必要时再粘贴 User-Agent。

也可以直接在命令行传入：

\begin{lstlisting}
./submit login --cookie 'UOJSESSID=...; cf_clearance=...' \
  --user-agent 'Mozilla/5.0 ...'
\end{lstlisting}

其它参数：

\begin{center}
\begin{tabularx}{\textwidth}{lX}
  \toprule
  参数 & 说明 \\
  \midrule
  \opt{--site} & OJ 地址，默认 \url{https://qoj.ac} \\
  \opt{--cookie} & Cookie 字符串 \\
  \opt{--cookie-file} & 从文件读取 Cookie \\
  \opt{--user-agent} & 浏览器 User-Agent \\
  \opt{-u}/\opt{--username} & 用户名（备选，密码登录） \\
  \opt{-p}/\opt{--password} & 密码（备选，不推荐） \\
  \bottomrule
\end{tabularx}
\end{center}

配置写入：

\begin{lstlisting}
~/.config/qoj-xcpc/config.json
\end{lstlisting}

\noindent Windows 上对应 \path{\%APPDATA\%/qoj-xcpc/config.json}。该文件含会话凭证，\textbf{不要}提交到 git。

\subsection{检查登录状态}

\begin{lstlisting}
./submit whoami
\end{lstlisting}

成功时会打印当前用户名与站点。Cookie 过期后重新执行 \cmd{./submit login}。

\section{绑定比赛}

每场 VP 在独立目录中进行。先在网页打开比赛；若尚未开始虚拟参赛，可在网页点 \textbf{Start}，或在 \cmd{init} 时加上 \opt{--start}。

\begin{lstlisting}
mkdir my-vp && cd my-vp
/path/to/submit init https://qoj.ac/contest/1234 --template --start
\end{lstlisting}

\cmd{init} 接受完整比赛 URL，或纯数字比赛 ID（此时使用已保存的站点，默认 \url{https://qoj.ac}）。\url{https://training.qoj.ac}（Public Judge）同样支持，请给出完整 URL。

\subsection{init 会做什么}

\begin{itemize}
  \item 解析比赛题目列表，写入当前目录的 \cmd{.qoj.json}。
  \item 将 \cmd{submit} 复制到当前目录（若尚不在此）。
  \item \opt{--template}：为每道题生成空模板 \cmd{a.cpp}、\cmd{b.cpp}、$\ldots$（已存在的文件不会覆盖）。
  \item \opt{--start}：尝试自动点击 VP 的 Start。
\end{itemize}

绑定成功后，目录大致如下：

\begin{lstlisting}
my-vp/
  submit          # 本插件
  .qoj.json       # 本场题目列表
  a.cpp
  b.cpp
  ...
\end{lstlisting}

之后所有提交相关命令都应在该目录（或其子目录向上能找到 \cmd{.qoj.json} 的位置）执行。

若 \cmd{init} 未能自动开始 VP，可再执行：

\begin{lstlisting}
./submit start
\end{lstlisting}

\section{比赛中提交}

这是插件的核心用法，与 XCPC 现场一致：\textbf{文件名第一个字母即题号}。

\begin{lstlisting}
./submit a.cpp
./submit b.cpp
./submit c.py
\end{lstlisting}

也可省略扩展名，脚本会按 \cmd{.cpp}、\cmd{.cc}、\cmd{.py} 等顺序查找：

\begin{lstlisting}
./submit a
\end{lstlisting}

提交时会打印比赛名、题号、语言与 URL，然后上传代码并轮询评测结果。按 \cmd{Ctrl+C} 可停止等待（代码已经交上）。

\subsection{常用选项}

\begin{center}
\begin{tabularx}{\textwidth}{lX}
  \toprule
  选项 & 说明 \\
  \midrule
  \opt{-p A} / \opt{--problem A} & 文件名不是题号时，显式指定题目 \\
  \opt{-l C++23} / \opt{--lang C++23} & 覆盖自动选择的语言 \\
  \opt{-n} / \opt{--no-wait} & 提交后不等待评测 \\
  \opt{--wait 180} & 等待评测的超时秒数，默认 180 \\
  \opt{--dry-run} & 只预览提交信息，不真正提交 \\
  \opt{--debug} & 打印 HTTP 请求，便于排查 \\
  \bottomrule
\end{tabularx}
\end{center}

示例：

\begin{lstlisting}
./submit b.cpp --lang C++23
./submit sol.cpp -p D
./submit a.cpp --dry-run
./submit a.cpp -n
\end{lstlisting}

\subsection{题号规则}

题号取自源文件\textbf{主文件名的第一个字母}（大小写不敏感）：

\begin{center}
\begin{tabular}{ll}
  \toprule
  文件 & 题目 \\
  \midrule
  \cmd{a.cpp}、\cmd{A.cpp} & A \\
  \cmd{b.cc} & B \\
  \cmd{c.py} & C \\
  \cmd{d.java} & D \\
  \bottomrule
\end{tabular}
\end{center}

若文件名无法对应本场题目，使用 \opt{-p}。本场题目可用 \cmd{./submit problems} 查看。

\subsection{默认语言}

脚本读取题目页上的语言列表，按扩展名优先匹配。题目不提供首选语言时，会沿列表向后兼容。

\begin{center}
\begin{tabular}{lll}
  \toprule
  扩展名 & 优先语言 & 其后回退 \\
  \midrule
  \cmd{.cpp} \cmd{.cc} \cmd{.cxx} & C++20 & C++23, C++17, $\ldots$ \\
  \cmd{.c} & C11 & C99 \\
  \cmd{.py} & PyPy3 & CPython 3 / Python 3 \\
  \cmd{.java} & Java 21 & Java 17, 11, 8 \\
  \cmd{.rs} & Rust & \\
  \cmd{.go} & Go & \\
  \cmd{.kt} & Kotlin & \\
  \cmd{.pas} & Pascal & \\
  \cmd{.d} & D & \\
  \cmd{.rb} & Ruby & \\
  \cmd{.hs} & Haskell & \\
  \bottomrule
\end{tabular}
\end{center}

\section{其它命令}

\begin{center}
\begin{tabularx}{\textwidth}{lX}
  \toprule
  命令 & 说明 \\
  \midrule
  \cmd{./submit problems} & 列出本场题目（别名 \cmd{list}） \\
  \cmd{./submit status} & 最近提交；\opt{-n 15} 控制条数 \\
  \cmd{./submit watch 123456} & 跟踪指定提交号的评测 \\
  \cmd{./submit start} & 开始 Virtual Participation \\
  \cmd{./submit whoami} & 检查当前登录账号 \\
  \bottomrule
\end{tabularx}
\end{center}

\section{配置文件}

\subsection{全局配置}

路径：\texttt{\textasciitilde/.config/qoj-xcpc/config.json}。主要字段：

\begin{lstlisting}
{
  "site": "https://qoj.ac",
  "cookie": "UOJSESSID=...; cf_clearance=...",
  "user_agent": "Mozilla/5.0 ...",
  "username": "your_handle"
}
\end{lstlisting}

\subsection{比赛配置}

路径：从当前目录向上查找的 \cmd{.qoj.json}（由 \cmd{init} 生成）。主要字段：

\begin{lstlisting}
{
  "site": "https://qoj.ac",
  "contest_id": "1234",
  "contest_url": "https://qoj.ac/contest/1234",
  "name": "Contest Name",
  "username": "your_handle",
  "problems": {
    "A": {"id": "3835", "title": "...", "url": "..."},
    "B": {"id": "3836", "title": "...", "url": "..."}
  }
}
\end{lstlisting}

\section{常见问题}

\begin{description}
  \item[被 Cloudflare 拦截]
    Cookie 必须与 User-Agent 来自同一浏览器会话，通常要带 \cmd{cf\_clearance}。过期后重新 \cmd{./submit login}。

  \item[解析不到题目]
    确认已经 Start VP、Cookie 有效，且 URL 形如 \url{https://qoj.ac/contest/<id>}。

  \item[提示尚未登录]
    先 \cmd{./submit login}，再用 \cmd{./submit whoami} 验证。

  \item[当前目录不是比赛目录]
    先 \cmd{./submit init <比赛URL>}，或 \cmd{cd} 到含有 \cmd{.qoj.json} 的目录。

  \item[交互题（Anna / Bruno）]
    插件会拒绝提交，请到网页上手交。

  \item[Public Judge]
    \url{https://training.qoj.ac} 同样支持，\cmd{init} 时使用完整 URL。

  \item[提交后一直 Waiting]
    可 \cmd{Ctrl+C} 退出等待，随后 \cmd{./submit status} 或打开提交页查看。
\end{description}

\section{命令速查}

\begin{lstlisting}
# 登录与检查
./submit login
./submit whoami

# 绑定比赛
./submit init https://qoj.ac/contest/1234 --template --start
./submit start
./submit problems

# 提交（XCPC 风格）
./submit a.cpp
./submit b.cpp --lang C++23
./submit sol.cpp -p D
./submit a.cpp --dry-run
./submit a.cpp -n

# 查询
./submit status
./submit watch 123456
\end{lstlisting}

\appendix
\section{编译本文}

\begin{lstlisting}
xelatex manual.tex
\end{lstlisting}

需要 CTeX 宏集（TeX~Live 或 MiKTeX 中的 \cmd{ctex} 包）。若中文无法显示，请确认已安装 XeLaTeX 与中文字体。

\end{document}
